% Generated draft: requires mathematical and source review.
\documentclass[11pt]{article}
\usepackage[margin=25mm]{geometry}
\usepackage{fontspec}
\setmainfont{Latin Modern Roman}
\usepackage{amsmath,amssymb,mathtools,unicode-math}
\setmathfont{Latin Modern Math}
\usepackage{enumitem,needspace,xurl,hyperref}
\setlength{\parindent}{0pt}
\setlength{\parskip}{5pt}
\setlength{\emergencystretch}{3em}
\newcommand{\N}{\mathbb{N}}
\newcommand{\Z}{\mathbb{Z}}
\newcommand{\Q}{\mathbb{Q}}
\newcommand{\R}{\mathbb{R}}
\newcommand{\C}{\mathbb{C}}
\newcommand{\F}{\mathbb{F}}
\newcommand{\A}{\mathbb{A}}
\newcommand{\PP}{\mathbb{P}}
\newcommand{\E}{\mathbb{E}}
\newcommand{\eps}{\varepsilon}
\newcommand{\dd}{\,\mathrm{d}}
\begin{document}
\section*{UnsolvedMath ID: 4600005}
\textbf{Source problem code:} catalogue entry
\subsection{Definitions and mathematical statement}
Let \(S\subseteq\mathbb R\) be a subring containing \(1\), and let \(\Lambda=(\lambda_1,\ldots,\lambda_k)\), \(k\ge1\), be a finite multiset of nonzero complex numbers, with multiplicities. A matrix \(A\) over \(S\) is primitive if all its entries are nonnegative and some positive integer power \(A^m\) has every entry strictly positive. Realizing \(\Lambda\) as its nonzero spectrum means that, for some size \(d\ge k\),
\[
 A\in M_d(S\cap\mathbb R_{\ge0}),\qquad A^m>0\text{ entrywise for some }m\ge1,
 \qquad
 \det(tI_d-A)=t^{d-k}\prod_{i=1}^k(t-\lambda_i).
\]
The number of additional zero eigenvalues is unrestricted.

The source's spectral conjecture asks for equivalence between this realization property and the ``Perron, Galois, and nonnegative net-trace conditions stated in the source''. Those conditions are not actually included in the supplied excerpt. In symbolic form the requested equivalence is
\[
 \Lambda\text{ has the realization above over }S
 \quad\Longleftrightarrow\quad
 \Lambda\text{ satisfies the source's three conditions for }S.
\]
It cannot be completed as a precise necessary-and-sufficient theorem without those missing conditions.

For orientation only, the usual Perron condition involves a positive real member strictly larger in modulus than all other members; the characteristic-polynomial identity also forces \(\prod_i(t-\lambda_i)\in S[t]\). The power traces are \(p_j(\Lambda)=\sum_i\lambda_i^j\). Expressions commonly called net traces involve sums \(\sum_{d\mid j}\mu(d)p_{j/d}(\Lambda)\), with \(\mu\) the number-theoretic M\"obius function. These observations do not identify the exact source conditions, their strictness or integrality clauses, or which trace requirements apply to which ring. In particular an integer-matrix net-trace criterion must not be imposed on all real coefficient rings without the original statement.

The registry labels the conjecture solved. Its summary reports a literature solution for \(S=\mathbb Z,\mathbb Q,\mathbb R\), calls these the main cases, and says the arbitrary-unital-subring formulation is subsumed or is not the substantive open core. No theorem citation or reduction for other subrings is supplied. Results for those three rings alone do not specify realization over every intermediate subring, so the original universal domain is retained and its reported general status remains unverified. The historical 2008 open label has no recorded later authoritative resolution check.
\subsection{Short English statement}
Characterize which finite multisets can be the nonzero spectrum of a primitive nonnegative matrix over a given unital subring of \(\mathbb R\), allowing extra zero eigenvalues. The supplied record omits the proposed exact conditions and provides no citations for its solved-status report, so the equivalence needs the original source.
\subsection{Sources}
\begin{itemize}
\item Local registry: \url{data/unsolvedmath/problems.json}, permanent UnsolvedMath ID \(4600005\), source code catalogue entry; supplied statement, background, and research summary, updated 2026-08-06.
\item Boyle, \emph{Open Problems in Symbolic Dynamics} (2008), Problem/Question/Conjecture 6.1: \url{https://www.math.umd.edu/~mboyle/papers/openfinalsub3nov2007.pdf}. Supplied PDF-text extraction accessed 2026-07-29; no later authoritative resolution check is recorded.
\item The summary asserts solutions for \(\mathbb Z,\mathbb Q,\mathbb R\) and a relation to arbitrary subrings, but supplies no identifiable solution publications or exact theorem statements.
\end{itemize}
\end{document}
